% unidades/29-keigo-sonkeigo.tex
\chapter{🙇 \ruby{尊敬語}{そんけいご} — Keigo: Verbos de Respeto}
\unidadheader{29}{尊敬語}{いらっしゃる・おっしゃる・なさる}{Expresar respeto elevando las acciones del interlocutor}

\begin{tcolorbox}[vocabbox, title={\emoji{pushpin} Vocabulario}]
\begin{tabularx}{\linewidth}{llX}
\toprule
\textbf{Japonés} & \textbf{Español} & \textbf{Tipo} \\ \midrule
\vocabitem{尊敬語}{そんけいご}{lenguaje de respeto (keigo)}{sustantivo}
\vocabitem{上司}{じょうし}{superior / jefe}{sustantivo}
\vocabitem{先生}{せんせい}{maestro / doctor}{sustantivo}
\vocabitem{お客様}{おきゃくさま}{cliente (muy formal)}{sustantivo}
\vocabitem{ご覧になる}{ごらんになる}{ver (honorífico)}{v. G1 irr.}
\vocabitem{ご存知}{ごぞんじ}{saber (honorífico)}{adj. na}
\vocabitem{めしあがる}{めしあがる}{comer/beber (honorífico)}{v. G1}
\bottomrule
\end{tabularx}
\end{tcolorbox}

\subsection*{\emoji{speech-balloon} Frases Clave}
\frase{\ruby{先生}{せんせい}はどちらにいらっしゃいますか？}{¿Dónde está usted, profesor?}
\frase{\ruby{部長}{ぶちょう}は\ruby{何}{なん}とおっしゃいましたか？}{¿Qué dijo el jefe de sección?}
\frase{\ruby{社長}{しゃちょう}はすでにご\ruby{覧}{らん}になりましたか？}{¿Ya lo vio el director?}
\frase{ご\ruby{存知}{ぞんじ}でしたか？}{¿Lo sabía usted?}

\subsection*{\emoji{gear} Gramática}
\patron{いらっしゃる — «estar / ir / venir» (honorífico)}
  {います / います / きます → いらっしゃいます}
  {\ej{\ruby{先生}{せんせい}は\ruby{今}{いま}\ruby{研究室}{けんきゅうしつ}にいらっしゃいます。}{El profesor está ahora en el laboratorio.}
   \ej{\ruby{社長}{しゃちょう}は\ruby{明日}{あした}こちらにいらっしゃいます。}{El director vendrá aquí mañana.}
   \ej{どちらからいらっしゃいましたか？}{¿De dónde vino usted?}}

\patron{おっしゃる / なさる / めしあがる — verbos honoríficos especiales}
  {いう → おっしゃる \quad する → なさる \quad たべる/のむ → めしあがる}
  {\ej{\ruby{先生}{せんせい}は\ruby{何}{なん}とおっしゃいましたか？}{¿Qué dijo el profesor?}
   \ej{\ruby{部長}{ぶちょう}はどんなお\ruby{仕事}{しごと}をなさっていますか？}{¿Qué tipo de trabajo hace el jefe de sección?}
   \ej{もう\ruby{昼食}{ちゅうしょく}はめしあがりましたか？}{¿Ya almorzó usted?}}

\patron{ご覧になる / ご存知 / お〜になる — patrón general}
  {みる → ご\ruby{覧}{らん}になる \quad しる → ご\ruby{存知}{ぞんじ} \quad V-stem: お〜になる}
  {\ej{この\ruby{映画}{えいが}はもうご\ruby{覧}{らん}になりましたか？}{¿Ya vio esta película?}
   \ej{\ruby{田中}{たなか}\ruby{部長}{ぶちょう}のことはご\ruby{存知}{ぞんじ}ですか？}{¿Conoce al jefe Tanaka?}
   \ej{\ruby{先生}{せんせい}はもうお\ruby{帰}{かえ}りになりましたか？}{¿El profesor ya se fue a casa? (お帰りになる)}}

\comparacion{尊敬語 vs 謙譲語}
  {尊敬語\\\small eleva las acciones del OTRO\\\small \ruby{先生}{せんせい}がいらっしゃる\\\textit{(el profesor «está» — honorífico)}}
  {謙譲語\\\small baja las acciones PROPIAS\\\small \ruby{私}{わたし}が\ruby{参}{まい}る\\\textit{(yo «voy» — humilde)}}
  {Regla fundamental: 尊敬語 describe lo que hacen los demás (superiores, clientes); 謙譲語 describe lo que hace uno mismo ante los superiores. Nunca use 尊敬語 para hablar de sus propias acciones, ni 謙譲語 para hablar de las acciones de los demás.}

\coloquial{〜いらっしゃいますか vs います (contextos mixtos)}{
  En situaciones semiformales, se puede mezclar niveles: con clientes en tiendas se usa 尊敬語, pero entre colegas del mismo nivel se usa 普通語. La clave es leer el contexto. El uso incorrecto de keigo (muy pomposo o muy casual) crea más incomodidad que no usarlo. Mejor seguro: いらっしゃいますか con personas mayores o clientes.}

\culturabox{\ruby{敬語}{けいご}と\ruby{社会構造}{しゃかいこうぞう} — Keigo y la jerarquía japonesa}{
  El keigo no es solo gramática: es un sistema que refleja y reproduce la estructura social japonesa. 上下関係 (jerarquía vertical) y 内外関係 (dentro/fuera del grupo) determinan qué formas usar. En empresas, se eleva al cliente siempre (お客様 = rey), se eleva al superior externo, y se humilla el hablante. Esta danza lingüística es una habilidad social crucial en el mundo laboral japonés.
}

\lectura{\ruby{電話対応}{でんわたいおう}\\[0.4em]
  「はい、ABC\ruby{商事}{しょうじ}でございます。」「\ruby{山田}{やまだ}\ruby{部長}{ぶちょう}はいらっしゃいますか？」「\ruby{申}{もう}し\ruby{訳}{わけ}ございません。\ruby{山田}{やまだ}はただいま\ruby{外出}{がいしゅつ}しております。」「\ruby{何時頃}{なんじごろ}お\ruby{戻}{もど}りになりますか？」「3\ruby{時}{じ}\ruby{頃}{ごろ}には\ruby{戻}{もど}る\ruby{予定}{よてい}でございます。よろしければ、ご\ruby{伝言}{でんごん}を\ruby{承}{うけたまわ}りましょうか？」}
{Atención telefónica\\[0.4em]
  «Sí, ABC Shoji, a sus órdenes.» «¿Está el jefe Yamada?» «Lo siento mucho. Yamada está fuera en este momento.» «¿A qué hora regresa?» «Tiene previsto regresar alrededor de las 3. Si lo desea, ¿puedo tomar un mensaje?»}

\begin{tcolorbox}[ejerciciobox, title={\emoji{pencil} Ejercicios}]
\begin{enumerate}
  \item Convierte a 尊敬語: 「先生はいますか？」「先生は何を食べましたか？」
  \item ¿Cuál es el honorífico de する? ¿Y de 言う? ¿Y de 見る?
  \item Crea un diálogo breve con al menos 2 formas 尊敬語 distintas.
\end{enumerate}
\textbf{Respuestas:}
1) 「\ruby{先生}{せんせい}はいらっしゃいますか？」「\ruby{先生}{せんせい}は\ruby{何}{なに}をめしあがりましたか？」 \quad
2) する→なさる、\ruby{言}{い}う→おっしゃる、\ruby{見}{み}る→ご\ruby{覧}{らん}になる \quad
3) (libre)
\end{tcolorbox}

\begin{tcolorbox}[kanjibox, title={\emoji{mahjong} Kanjis}]
\begin{tabularx}{\linewidth}{cllXX}
\toprule
\textbf{Kanji} & \textbf{On} & \textbf{Kun} & \textbf{Significado} & \textbf{Vocab} \\ \midrule
\kanjirow{尊}{そん}{\ruby{たっと}{い}}{venerable/respetar}{\ruby{尊敬}{そんけい} / \ruby{尊重}{そんちょう}}
\kanjirow{敬}{けい}{\ruby{うやま}{う}}{respetar}{\ruby{尊敬}{そんけい} / \ruby{敬語}{けいご}}
\kanjirow{覧}{らん}{—}{ver/contemplar}{\ruby{ご覧}{ごらん} / \ruby{一覧}{いちらん}}
\kanjirow{存}{そん/ぞん}{—}{existir/saber}{\ruby{存知}{ぞんじ} / \ruby{存在}{そんざい}}
\bottomrule
\end{tabularx}
\end{tcolorbox}
